\section{功能安全与确定性硬件：在故障容忍时间内达到受控状态}

可靠性关注部件或系统是否持续正确工作，功能安全关注故障发生后是否会产生不可接受的危险。一个控制器可以因检测到内部错误而主动停机：从可用性看它中断了服务，从功能安全看它可能及时阻止了错误执行器输出。反过来，系统仍然响应网络请求并不表示安全；若传感器已经失真而输出仍持续更新，服务“在线”可能掩盖危险状态。IEC 61508 给出电气、电子与可编程电子安全相关系统的通用框架，ISO 26262 将功能安全生命周期落实到道路车辆电子电气系统及硬件开发。\cite{iec61508}\cite{iso26262-hardware}

\topic{安全目标先定义危险输出，再给检测与反应分配时间}

\term{安全目标}{safety goal}不是“系统不能出错”，而是对危险行为给出可验证限制。例如，传感器无效时执行器不得继续增加输出；主控制器失去可信状态后，独立路径必须在规定时间内关闭驱动；通信丢失后系统只能保持有限输出一段确定时间。每个目标都要明确受保护对象、触发条件、允许的降级行为和最终安全状态。

\term{故障容忍时间间隔}{Fault Tolerant Time Interval, FTTI}是从故障发生到危险后果必须被阻止之间可使用的最大窗口。检测、确认、仲裁、通信、执行器动作和物理惯性都占用该窗口。平均响应时间不能证明安全，因为一次超过 FTTI 的尾部事件就可能使安全机制失效。

\begin{pdfdiagram}[x=1cm,y=1cm]
  \node[hw block, minimum width=2.05cm, align=center] (fault) at (0.7,2.9)
    {故障出现\\传感器卡滞};
  \node[hw state, minimum width=2.05cm, align=center] (observe) at (3.2,2.9)
    {取得证据\\时间戳与双通道};
  \node[hw compute, minimum width=2.05cm, align=center] (decide) at (5.7,2.9)
    {安全判定\\新鲜度与一致性};
  \node[hw control, minimum width=2.05cm, align=center] (act) at (8.2,2.9)
    {输出门控\\撤销驱动许可};
  \node[hw state, minimum width=2.05cm, align=center] (safe) at (10.7,2.9)
    {受控状态\\执行器已关闭};
  \draw[hw bus] (fault) -- (observe) -- (decide) -- (act) -- (safe);

  \draw[draw=black!65, line width=0.7pt] (0.7,0.75) -- (10.7,0.75);
  \draw[draw=black!65, line width=0.7pt] (0.7,0.55) -- (0.7,0.95);
  \draw[draw=black!65, line width=0.7pt] (10.7,0.55) -- (10.7,0.95);
  \node[hw label] at (5.7,0.35) {FTTI：从故障到阻止危险结果的总预算};
  \draw[hw danger] (safe.south) -- node[right,hw label]{超过此点即不能证明安全} (10.7,1.55);
\end{pdfdiagram}
\diagramnote{FTTI 是端到端时间合同。故障检测只是中间状态；只有执行器输出已经被硬件门控到安全值，或系统已进入经分析允许的降级模式，安全动作才完成。日志写入和维修通知可以随后进行，但不能阻塞关键关断路径。}

假设某执行器安全目标给出 10 ms FTTI。设计预算可能包括：采样数据最老 1.5 ms，输入滤波与双通道比较 0.5 ms，安全任务最坏执行 2.0 ms，消息与门控传播 1.0 ms，执行器电气和机械撤销 3.0 ms，故障确认与状态切换 1.0 ms。合计 9.0 ms，只剩 1.0 ms 集成裕量。若软件任务平均只需 0.4 ms，却偶尔因 Cache miss、总线拥塞或不可屏蔽临界区达到 3 ms，平均值对该安全目标没有证明力。

\begin{longtable}{L{0.19\linewidth}L{0.24\linewidth}L{0.28\linewidth}L{0.19\linewidth}}
\toprule
预算对象 & 需要约束的状态 & 适合的硬件措施 & 验收证据 \\
\midrule
输入年龄 & 采样到消费的最大间隔 & 硬件时间戳、双缓冲与过期标志 & 最旧允许样本仍满足预算 \\\tablerowrule
计算时延 & 最坏路径、抢占与存储访问 & TCM、确定性 Cache 策略、优先级和预算计数器 & WCET 分析与高压干扰测试 \\\tablerowrule
通信时延 & 仲裁、重试和拥塞上限 & 保留时隙、优先级、credit 与独立安全通道 & 最大负载下的端到端上限 \\\tablerowrule
关断传播 & 门控、驱动器和物理惯性 & 硬件许可链、独立电源开关与反馈触点 & 命令到物理安全状态的实测时间 \\\tablerowrule
恢复时间 & 诊断、重启和重新许可 & 安全状态锁存、分阶段自检与双通道确认 & 未通过验收前输出保持禁止 \\
\bottomrule
\end{longtable}

\topic{确定性不是速度高，而是最坏边界可证明}

实时处理器常提供紧耦合存储器（TCM）、局部中断控制器、受控 MPU、可预测流水线和明确的错误响应。TCM 牺牲自动替换带来的平均性能，换取固定地址范围内不发生 Cache miss；静态分配的总线时隙或高优先级虚通道限制仲裁等待；中断优先级、嵌套深度和最大屏蔽时间共同决定入口上界。Arm Cortex-R 类处理器把确定性、隔离和功能安全作为同一类系统需求处理，并支持锁步等诊断组织。\cite{arm-lockstep-safety}

Cache、乱序执行和推测并非绝对不能用于安全任务，但必须给出可验证策略。可以在任务释放前预装并锁定关键代码/数据，划分 Cache way，限制 DMA 与其他核心的共享带宽，或把安全关键路径放进 TCM。若最坏行为取决于另一个不受控租户是否产生大量 miss，就不能只用一次基准测试声明满足截止时间。

\begin{longtable}{L{0.18\linewidth}L{0.25\linewidth}L{0.29\linewidth}L{0.18\linewidth}}
\toprule
共享资源 & 不确定性来源 & 形成上界的方法 & 不能替代证明的指标 \\
\midrule
Cache/TLB & 替换、冲突、页表遍历与 shootdown & 锁定、分区、预热、TCM 或计入最坏 miss & 平均命中率 \\\tablerowrule
DRAM/NoC & bank 冲突、刷新、仲裁和其他主设备 & 带宽保留、优先级上限、时隙与干扰模型 & 空载带宽 \\\tablerowrule
中断 & 屏蔽区、嵌套、入口保存和中断风暴 & 最大屏蔽时间、优先级规则和速率限制 & 平均中断延迟 \\\tablerowrule
DMA/设备 & burst、重试、完成聚合和链路恢复 & 描述符长度上限、独立队列与超时合同 & 峰值设备吞吐 \\\tablerowrule
时钟/DVFS & 频率切换、热降频和时钟漂移 & 安全档固定频率、状态确认和最坏降频预算 & 标称最高频率 \\
\bottomrule
\end{longtable}

\topic{锁步检测分歧，安全岛决定分歧之后做什么}

\term{双核锁步}{Dual-Core Lockstep, DCLS}让两个处理器核心接收等价输入并执行同一指令流，在架构输出或接口边界比较结果。比较器检测到分歧时，不能据此判断哪个核心正确；它只能证明至少一个结果不可信。因此锁步后还需要独立安全控制器保存故障位置、阻断主输出、选择受控降级或进入安全状态。

两个副本应尽量减少共同原因。延迟锁步可以让第二核心相对第一核心错开若干周期，使一次窄电压或时钟扰动不同时命中同一内部阶段；独立电源监视、时钟监视和物理布局分离可进一步降低共同故障概率。共享 Cache、互连和电源仍要使用 ECC、奇偶校验、CRC、超时或独立监视器保护。只复制 CPU，而把比较器、复位线和执行器许可都放在同一无保护控制块中，会把故障集中到新的单点。

\begin{pdfdiagram}[x=1cm,y=1cm]
  \node[hw state, minimum width=2.55cm, align=center] (input) at (0.8,3.55)
    {受保护输入\\时间戳、CRC、新鲜度};
  \node[hw compute, minimum width=2.55cm, align=center] (corea) at (4.2,4.45)
    {核心 A\\安全控制任务};
  \node[hw compute, minimum width=2.55cm, align=center] (coreb) at (4.2,2.65)
    {核心 B\\延迟锁步副本};
  \node[hw control, minimum width=2.55cm, align=center] (cmp) at (7.6,3.55)
    {锁步比较器\\接口状态逐拍核对};
  \node[hw control, minimum width=2.75cm, align=center] (island) at (10.9,3.55)
    {安全岛\\分类、门控、记录};

  \node[hw block, minimum width=2.55cm, align=center] (monitor) at (4.2,0.55)
    {独立监视器\\电压、时钟、watchdog};
  \node[hw state, minimum width=2.75cm, align=center] (gate) at (10.9,0.55)
    {执行器许可门\\默认关闭};

  \draw[hw bus] (input) -- (corea);
  \draw[hw bus] (input) -- (coreb);
  \draw[hw bus] (corea) -- (cmp);
  \draw[hw bus] (coreb) -- (cmp);
  \draw[hw bus] (cmp) -- (island);
  \draw[hw danger] (monitor) -- (island);
  \draw[hw bus] (island) -- (gate);
  \draw[hw return] (gate.west) -- node[hw label,below]{反馈确认已到安全值} (monitor.east);
\end{pdfdiagram}
\diagramnote{锁步系统的责任边界。两个核心只提供分歧证据；安全岛还接收独立电压、时钟和 watchdog 证据，并控制默认关闭的执行器许可门。许可门反馈确认后，故障反应才完成。}

\topic{诊断覆盖率必须包含安全机制自身}

诊断覆盖率可以用“目标故障集合中能在规定时间内被检测并正确响应的比例”理解。分母必须来自故障模型：瞬态 bit 翻转、永久 stuck-at、时钟停止、总线超时、地址译码错误、传感器开路或短路等不能混成一个无边界数字。检测到错误却来不及阻止危险输出，不计入该安全目标的有效覆盖。

安全机制也会失效。watchdog 计数器可能卡住，锁步比较器可能恒等，ECC 注入路径可能没有真正穿过译码器，执行器门控晶体管也可能短路。平台要周期性测试这些监视与关断通路：制造测试验证结构覆盖，启动 BIST 建立初始状态，运行期测试在不产生危险输出的窗口轮换检查潜伏故障。双通道系统若长期只使用通道 A，通道 B 的失效可能一直潜伏到真正需要切换时。

\begin{longtable}{L{0.17\linewidth}L{0.24\linewidth}L{0.27\linewidth}L{0.22\linewidth}}
\toprule
故障类别 & 首个可靠证据 & 安全动作 & 禁止推断 \\
\midrule
核心分歧 & 锁步比较器、控制流签名或接口不一致 & 阻断输出，保存比较点，交安全岛处理 & 任选一个核心继续即正确 \\\tablerowrule
存储器错误 & ECC syndrome、地址和累计趋势 & 纠正后记录；不可纠正时隔离与降级 & ECC 已纠正便无需维护 \\\tablerowrule
时钟/电源越界 & 独立监视器与保持时间窗口 & 冻结提交、撤销许可、进入安全档 & CPU 仍执行便说明时序可靠 \\\tablerowrule
通信超时 & sequence、时间戳和最后有效消息年龄 & 使用限定时间的降级值，再转安全状态 & 旧消息格式正确便仍可使用 \\\tablerowrule
安全机制潜伏故障 & 周期自检或故障注入未得到预期反应 & 下线通道、限制功能并安排维修 & 未触发告警等于机制健康 \\
\bottomrule
\end{longtable}

\topic{恢复不是复位完成，而是重新建立安全许可}

功能安全状态机通常至少区分 Normal、Degraded、Safe 和 Service。Normal 允许全部功能；Degraded 在已分析的能力与时限内继续有限服务；Safe 撤销危险输出并保持最小受控状态；Service 允许诊断，但不能自动恢复任务期许可。错误消失不应直接从 Safe 跳回 Normal，因为传感器、输出门、通信序号和安全机制可能仍保留旧状态。

恢复路径先确认危险能量已被隔离，再执行自检、校准和双通道一致性检查，重建时间基准与通信 generation，最后由独立条件共同授予输出许可。若系统在 FTTI 内只能到达 Degraded，则安全分析必须证明该模式在限定时间内仍不会产生危险，并给出继续转入 Safe 的第二个截止时间。这样，确定性执行、锁步、诊断覆盖和状态恢复才共同构成功能安全闭环。
